% ------------------------------------------------------------------------
% file `main-bookexercise-3-solution-body.tex'
%
%     solution of type `bookexercise' with id `3'
%
% generated by the `exerciseanswer' environment of the
%   `xsim' package v0.21 (2022/02/12)
% from source `main' on 2026/10/10 on line 704
% ------------------------------------------------------------------------
1. 选C。一个最快的方法就是判断节点是否存在环。根据树的定义，若存在环则前驱后继关系一定不符合树的定义。笔者不建议在此处逻辑上浪费过多时间。\\

2. 选C。根据度的定义知道A、B、D正确。二叉树中任何一个节点至多有两棵子树，C错误。\\

3. 选B。树的高度和树的深度一定相等，并且等于树中层次的数量。该图中的树有四层，因此树的高度与深度都是4，A正确；D显然正确。\\

4. 如下表。
\begin{center}
    \begin{tabular}{|c|c|c|}
        \hline
         & 度为 0 的节点数 & 度为 2 的节点数 \\
        \hline
        图一 & 2 & 1 \\
        \hline
        图二 & 3 & 2 \\
        \hline
        图三 & 3 & 2 \\
        \hline
        图四 & 4 & 3 \\
        \hline
    \end{tabular}
\end{center}
事实上，\textbf{记二叉树中度为2的节点的数量为\(n_2\), 度为0的节点（即叶节点）数为\(n_0\)， 则\(n_0\) = \(n_2\)+1 对任何非空二叉树均成立。}对于数学上严谨的证明，
感兴趣的同学可以参考相关资料，我们这里只给出直观上的理解。

相信你也已经看出来了，最初对于只有一个节点的二叉树，\(n_0\) = \(n_2\)+1显然成立。此后，在二叉树中增加节点，只可能从度为0的节点“引出边”，或者从度为1的
节点“引出边”。\\
1）从一个度为0的节点“引出一条边”，不会增加\(n_2\)的数量，并且\(n_0\)会被吃掉一个，也会新增一个，相当于没变。所以该操作不影响\(n_0\)和\(n_2\)的数量关系；\\
2）从一个度为1的节点“引出一条边”，则会将\(n_0\)和\(n_2\)的数量都增加一个。该操作同样不影响\(n_0\)和\(n_2\)的数量关系；

所以，只要二叉树中至少有一个节点，此后无论加入多少个节点，\(n_0\)和\(n_2\)的关系都不会变化。\(n_0\) = \(n_2\)+1对任何非空二叉树都成立。\textbf{事实上
这是二叉树中一个非常重要的性质}，很多教材
喜欢将此直接作为“二叉树的性质”给出，但笔者认为，探索、发现规律的过程可以锻炼能力、培养创造力和解决问题的能力。并且这对知识的记忆会是一个很强的冲击。如果
直接记住公式，非常容易忘，并且很可能出现“\(n_0\)是啥意思来着？是谁等于谁加一来着？”这样的现象。而有了这个过程，当你忘记的时候，拿出笔画几个圈，立即能回忆起来。\\

5. 由{\color{red}\textbf{上一题}}的结论可知，非空二叉树中度为2的节点\(n_2 = 8\)，则度为0的节点数一定是8 + 1 = 9，即叶节点一定是9个。题目又说了度为1的节点
有3个，所以二叉树中节点总数为\(n = n_0 + n_1 + n_2 = 9 + 3 + 8 = 20\)。\\

6. 由已知信息可知，有\(n_0 + n_2 = 13\)成立。又根据前面的结论有\(n_0 = n_2 + 1\)成立，解二元一次方程，得\(n_0 = 7\)。即叶节点有7个。\\

7. 是满二叉树的有\underline{C、F}；不是满二叉树，但是是完全二叉树的有\underline{A、D、E}。完全二叉树只能是在最底层、从右开始有连续的残缺。F项倒数
第二层存在残缺；B项最底层的残缺并非从右开始并且连续，因此F、B都不是完全二叉树。\\

8. 选A、B、D。对于此类问题，动手画一画图就非常容易得到答案。此处笔者提醒大家对D选项稍作留意。该选项是后面我们介绍\uwave{堆}时会用到的一个完全二叉树的性质。\\



